\subsection{Coordination should earn its complexity only where evidence is genuinely distributed}

A complementary cross-domain benchmark reports that coordinated scientific agents can improve inference when different channels each expose only part of the phenomenon, while in at least one task a strong combined-summary baseline is effectively tied with the coordinated workflow \cite{wong2026coordination}. That result is directly relevant to the matched Orion evaluation because it argues against treating decomposition or multi-agent coordination as an intrinsic advantage. The scientific question is conditional: under what evidence topology does the added control architecture improve a preregistered outcome that a simpler synthesis cannot match?

The Orion comparison therefore preserves strong simpler baselines rather than using a weak single-prompt straw man. It should stratify tasks by evidence distribution and failure mode where possible: single-dominant-signal tasks, genuinely partial cross-source evidence, contradiction-rich tasks, provenance-dependent tasks, and tasks whose main challenge is mechanism discrimination rather than information aggregation. If the benefit appears only in a subset, the result should be reported as a capability frontier rather than a global superiority claim. This makes ``simpler lawful parent wins'' a positive scientific outcome: it identifies the contexts in which the additional epistemic machinery is unnecessary.
